\chapter*{Preamble}
\addcontentsline{toc}{chapter}{Preamble}
\markboth{Preamble}{Preamble}

\articlenote{Read aloud at the naming of every child, at the seating of every Mayor, and at the opening of every school year, that none may say they did not know why they live as they live.}

The world above is poison, and the world below is ours to keep. This is the whole of what our forebears left us to remember, and it is enough. Beyond that single fact, this Pact does not ask any citizen to hold the old world in mind at all --- neither to mourn it nor to picture it nor to wonder after it --- for such wondering has never once led anywhere good, and the Pact exists, before it exists to do anything else, to keep us from wandering there.

We are told, and it is sufficient to be told, that the air outside once turned against the people who breathed it, and that in the turning, nearly everything that had been built was lost. What remained of the builders' foresight was this: a shaft sunk deep into the earth, sealed against the poison above, stocked and engineered to hold generations of the living long enough for the world to heal, if it ever will. We do not know how long that healing will take. We are not to guess. Guessing has its own poison, and it moves faster underground than any vapor moves above.

The silo is not a shelter merely. It is a vessel with a hull, and a hull cannot be pierced in one place without the whole of it taking on water. Every rule that follows in this Pact --- on birth, on labor, on speech, on punishment --- exists to keep that hull whole for as many turns of the generations as it takes. We ask each citizen to hold this thought above all others: that the silo does not merely contain your life, it is the only reason you have one to live, and every citizen after you depends on your keeping of these laws exactly as every citizen before you depended on the keeping of them by others, most of whom you will never know by name, all of whom kept the silo alive so that you might stand here reading this. A citizen who keeps this Pact in full is sound, in the same word and the same sense a wall or a seal is sound, and the silo's whole endurance is nothing more, and nothing less, than the sum of so many sound citizens as it can keep.

This is a Pact, and a Pact is a thing agreed to by more than one party. It is not handed down from a single hand, though it is set down by many hands together, to be amended in Assembly, argued over in Judicial, and tested against whatever winters of scarcity and years of grief the silo's generations must yet bear. It binds the Mayor no less than it binds the porter who carries her messages, the Head of Information Technology no less than the child not yet given a name. No officer of this silo, however seated, stands above it. Where any part of this Pact is silent, the silence is not permission; it is a gap to be carried to Judicial, and not filled by any one person's will.

We set down four founding truths, and every Article that follows serves one or more of them.

\textbf{First}, that the silo's people must go on, generation after generation, in numbers the silo can feed, warm, and air --- no more, for the silo does not forgive a burden it cannot carry, and no fewer, for a silo that dwindles is a silo that dies alone. This is the whole reason for the lottery and for every law of birth and household that follows from it.

\textbf{Second}, that the silo's work must go on --- the growing of food, the turning of the generator, the cleaning of the water, the keeping of the air --- and that this work is not beneath any citizen and is owed by every citizen, in the place the silo's own methods of testing and shadowing determine each is best suited to give it. No family, no floor, no trade stands higher than another in the sight of this Pact, whatever any resident may privately believe about the floor above or below.

\textbf{Third}, that the silo's peace must go on. We are too many, packed too closely, with too little room to run from one another, to survive our own quarrels unrestrained. Every law of conduct, of property, of marriage, of dispute, and of punishment set down here exists only to keep neighbor from turning on neighbor, and floor from turning on floor, in a hull that cannot afford the turning.

\textbf{Fourth}, and hardest, that the silo's mind must go on. The old world is gone, and the wanting of it is a sickness that has taken good men and women both, quietly at first and then all at once. This Pact does not pretend such wanting can be legislated out of the heart. It only insists that it be kept out of the mouth, off the walls, and away from the young, so that longing does not become teaching, and teaching does not become a whole floor's grief, and a whole floor's grief does not become the silo's end. Where this Pact is sternest, it is sternest here, and it does not apologize for that.

Let it also be said, for it is said too rarely: this Pact is not cruelty dressed as order. It is written by people who have buried their own dead and will raise their own children beneath its rule, and every hard clause in it is hard-won by someone's hard year, present or yet to come. It changes, by the lawful means set down in its final Article, when the silo's Assembly and its Mayor and its Judicial agree that it must, and it will change again, and again after that, for as long as the silo stands. What it does not do, and must never do, is disappear --- for a silo with no Pact is not a silo with more freedom, only a silo with no floor beneath its feet.

Keep it, then, not because you were told to, but because you have read, here, exactly why.
